\documentclass[12pt,a4paper]{article}
\usepackage[utf8]{inputenc}
\usepackage[italian]{babel}
\usepackage{hyperref}
\usepackage{graphicx}
\usepackage{geometry}
\usepackage{listings}
\usepackage{xcolor}

\geometry{a4paper, margin=2.5cm}

\definecolor{codegray}{rgb}{0.5,0.5,0.5}
\definecolor{codepurple}{rgb}{0.58,0,0.82}
\definecolor{backcolour}{rgb}{0.95,0.95,0.92}

\lstdefinestyle{mystyle}{
    backgroundcolor=\color{backcolour},   
    commentstyle=\color{codegray},
    keywordstyle=\color{magenta},
    numberstyle=\tiny\color{codegray},
    stringstyle=\color{codepurple},
    basicstyle=\ttfamily\footnotesize,
    breakatwhitespace=false,         
    breaklines=true,                 
    captionpos=b,                    
    keepspaces=true,                 
    numbers=left,                    
    numbersep=5pt,                  
    showspaces=false,                
    showstringspaces=false,
    showtabs=false,                  
    tabsize=2
}
\lstset{style=mystyle}

\title{\textbf{Manuale Completo e Dettagliato: Academic Knowledge Engine}}
\author{Antigravity Assistant}
\date{\today}

\begin{document}

\maketitle
\tableofcontents
\newpage

\section{Introduzione all'Academic Knowledge Engine}
Benvenuto nella guida completa del tuo \textbf{Academic Knowledge Engine}. Questo sistema non è un semplice "motore di ricerca" per file, ma un vero e proprio \textbf{Grafo della Conoscenza (Knowledge Graph)} combinato con intelligenza artificiale. 

L'obiettivo principale è fornire al tuo \textit{Pi Coding Agent} una memoria permanente e contestuale relativa ai tuoi studi universitari, permettendogli di aiutarti nello sviluppo software tenendo conto delle lezioni teoriche (es. Sicurezza, Reti, Sistemi Operativi) in modo preciso e affidabile, abbattendo le "allucinazioni" (ovvero le risposte inventate dall'IA).

Il sistema ingerisce i dati, li trasforma in frammenti (Chunk), genera delle rappresentazioni matematiche del significato (Embedding Vettoriali) e li collega tra loro su un database Neo4j.

\section{Architettura di Sistema e Requisiti}
Il sistema è costruito interamente in locale sul tuo Mac (Privacy-First) e si basa su:
\begin{itemize}
    \item \textbf{Python 3.12+}: Per l'elaborazione dei dati, l'estrazione dai PDF e l'interazione con le API.
    \item \textbf{Neo4j}: Il database a grafo che gestisce sia le connessioni logiche tra i concetti, sia gli indici vettoriali per la ricerca per similarità (RAG - Retrieval-Augmented Generation).
    \item \textbf{MCP (Model Context Protocol)}: L'interfaccia standard creata da Anthropic che permette al tuo agente IA di comunicare con il database.
\end{itemize}

\section{Avvio Iniziale e Setup Base}

\subsection{1. Accendere il Database Neo4j}
Neo4j deve essere sempre in esecuzione quando vuoi interrogare i dati o caricare nuovi file. Essendo installato nativamente tramite Homebrew, devi avviare il servizio:
\begin{lstlisting}[language=bash]
brew services start neo4j
\end{lstlisting}
Per fermarlo quando hai finito di lavorare (per risparmiare RAM):
\begin{lstlisting}[language=bash]
brew services stop neo4j
\end{lstlisting}

\subsection{2. Creare o Modificare la Password}
La prima volta, Neo4j richiede l'impostazione di una password per l'utente "neo4j". Se non lo hai già fatto o se ottieni un errore di autenticazione ("AuthError"), esegui questo comando da terminale:
\begin{lstlisting}[language=bash]
neo4j-admin dbms set-initial-password YOUR_PASSWORD_HERE
\end{lstlisting}

\subsection{3. Attivare l'Ambiente Python}
Apri il terminale, naviga nella cartella del progetto e attiva l'ambiente virtuale:
\begin{lstlisting}[language=bash]
cd /path/to/academic-kb
source .venv/bin/activate
\end{lstlisting}
Devi vedere \texttt{(.venv)} all'inizio della riga di comando.

\section{Configurazione Avanzata: Il file \texttt{.env}}
Nella cartella principale troverai un file chiamato \texttt{.env.example}. Rinominandolo o copiandolo in \texttt{.env} potrai configurare le chiavi.

\begin{lstlisting}[language=bash]
# Impostazioni Database
NEO4J_URI=bolt://localhost:7687
NEO4J_USERNAME=neo4j
NEO4J_PASSWORD=YOUR_PASSWORD_HERE

# Impostazioni Embeddings (Creazione Vettori)
EMBEDDING_PROVIDER=mock
EMBEDDING_MODEL=text-embedding-004
EMBEDDING_DIMENSION=768
# Se usi Gemini o OpenAI, metti qui la tua API Key:
EMBEDDING_API_KEY=tua_api_key_reale

# Impostazioni Notion
NOTION_API_KEY=secret_tua_chiave_notion
NOTION_ROOT_DATABASE_ID=tuo_id_database
\end{lstlisting}

\textbf{Importante:} Attualmente il sistema è impostato su \texttt{mock} (finto). Questo significa che genera vettori casuali perfetti per testare il caricamento offline senza consumare crediti. Quando vorrai che l'IA capisca *davvero* il significato delle parole, cambia \texttt{EMBEDDING\_PROVIDER} in \texttt{gemini} o \texttt{openai} e inserisci una chiave API valida.

\section{Upload e Analisi dei Dati (Ingestion Pipeline)}

\subsection{Caricamento dei PDF Universitari}
1. Crea una cartella in cui posizionare temporaneamente i tuoi file, ad esempio \texttt{cns\_pdfs}.
2. Trascina i file PDF (slide, dispense) al suo interno.
3. Lancia lo script python per farli digerire al sistema. Il parametro \texttt{--course} ti permette di taggare quei documenti in modo che appartengano a uno specifico corso universitario:
\begin{lstlisting}[language=bash]
python scripts/index.py --pdf-dir cns_pdfs --course "Computer Network Security"
\end{lstlisting}

\textbf{Cosa succede dietro le quinte?}
Il connettore PDF legge pagina per pagina. Invia il testo a un "Chunker" che lo taglia in paragrafi (es. 1000 caratteri) sovrapponendo leggermente le frasi finali per non tagliare concetti a metà (Sliding Window). Infine, calcola l'hash (checksum) del file per non ricaricarlo due volte se non è stato modificato.

\subsection{Collegare Notion (Sincronizzazione Incrementale)}
Vuoi che il sistema legga automaticamente i tuoi appunti presi su Notion? Segui questi passi millimetrici:
\begin{enumerate}
    \item Vai su \url{https://www.notion.so/my-integrations}.
    \item Clicca "New integration", dalle il nome "Academic KB" e seleziona il tuo workspace.
    \item Copia l'\textbf{Internal Integration Secret} e incollalo nel file \texttt{.env} sotto \texttt{NOTION\_API\_KEY}.
    \item Vai sulla pagina principale dei tuoi appunti su Notion. In alto a destra clicca i tre pallini (Menu) $\rightarrow$ \textbf{Connect to} $\rightarrow$ cerca la tua integrazione "Academic KB" e aggiungila.
    \item Copia il link della pagina Notion. L'ID del database è la sequenza di lettere e numeri lunghissima dopo l'ultimo \textit{slash} (/) e prima del punto interrogativo (\textit{?}). Incollalo in \texttt{NOTION\_ROOT\_DATABASE\_ID}.
    \item Avvia la sincronizzazione:
\end{enumerate}
\begin{lstlisting}[language=bash]
python scripts/index.py --notion
\end{lstlisting}
\textit{Il sistema confronterà la data di ultima modifica (\texttt{last\_edited\_time}) e aggiornerà solo i paragrafi che hai modificato o aggiunto, risparmiando tempo e chiamate API.}

\section{Come Funziona la Ricerca (Il Cuore dell'IA)}
Quando chiedi qualcosa a Pi Coding Agent, l'Academic Engine entra in gioco e compie una \textbf{Ricerca Ibrida (Hybrid Retrieval)}.

1. \textbf{Ricerca Semantica Vettoriale:} Converte la tua domanda ("Come fa il WEP ad autenticare?") in un vettore matematico. Cerca nello spazio i vettori più vicini. Questo trova i concetti per \textit{significato}, non per parole esatte.
2. \textbf{Ricerca Full-Text (BM25):} Molto spesso nel codice e nei concetti tecnici si cercano variabili esatte (\texttt{pthread\_mutex}). La ricerca semantica farebbe fatica, quindi entra in gioco la ricerca lessicale classica su Neo4j.
3. \textbf{Reciprocal Rank Fusion (RRF):} È l'algoritmo matematico che fonde i risultati di (1) e (2). Prende le liste di risultati, calcola un punteggio fuso basato sulla loro posizione ($Score = 1 / (60 + Rank)$) e restituisce la Top 5 dei frammenti più rilevanti e perfetti.

\section{Esplorazione Visiva del Grafo (Neo4j Browser)}
Il modo più potente per comprendere i dati è guardarli.
\begin{enumerate}
    \item Apri \url{http://localhost:7474} nel browser.
    \item Accedi (User: \texttt{neo4j}, Password: \texttt{academic\_secret\_password}).
    \item Nella console superiore, esegui comandi \textbf{Cypher} (il linguaggio di Neo4j). Ecco i più utili:
\end{enumerate}

\textbf{Per vedere tutti i documenti e come sono divisi in chunk:}
\begin{lstlisting}[language=SQL]
MATCH (d:Document)-[r:HAS_CHUNK]->(c:Chunk) RETURN d, c LIMIT 50
\end{lstlisting}

\textbf{Per contare quanti paragrafi hai per ogni corso universitario:}
\begin{lstlisting}[language=SQL]
MATCH (d:Document)-[:HAS_CHUNK]->(c:Chunk)
RETURN d.course as Corso, count(c) as Numero_Paragrafi
\end{lstlisting}

\textbf{Per cancellare tutto e ricominciare da zero (Attenzione!):}
\begin{lstlisting}[language=SQL]
MATCH (n) DETACH DELETE n
\end{lstlisting}

\section{Integrazione con Pi Coding Agent}
Per far parlare Pi con la tua base di conoscenza, devi aggiornare la sua configurazione MCP.
Nel file di configurazione del tuo client MCP (che si trova spesso nelle impostazioni dell'app o in un \texttt{mcp.json}), aggiungi il server:

\begin{lstlisting}[language=json]
{
  "mcpServers": {
    "academic-kb": {
      "command": "/path/to/academic-kb/.venv/bin/python",
      "args": ["/path/to/academic-kb/mcp_server/server.py"]
    }
  }
}
\end{lstlisting}

Fatto ciò, Pi avrà uno strumento chiamato \texttt{academic\_context}. Quando formulerai domande del tipo \textit{"Aiutami a scrivere questo codice C basandoti sulle lezioni di Sistemi Operativi"}, Pi eseguirà segretamente \texttt{academic\_context}, riceverà il testo fuso estratto dai tuoi PDF, e ti fornirà un codice accompagnato da citazioni reali dei tuoi professori!

\section{Risoluzione Problemi e Debugging}
\begin{itemize}
    \item \textbf{Errore di connessione (Errno 61):} Hai dimenticato di avviare il database con \texttt{brew services start neo4j}.
    \item \textbf{ImportError / ModuleNotFoundError:} Hai dimenticato di eseguire \texttt{source .venv/bin/activate} prima di eseguire \texttt{python scripts/index.py}.
    \item \textbf{L'IA non trova nulla di rilevante:} Hai lasciato \texttt{EMBEDDING\_PROVIDER=mock}. Questo crea vettori "finti". Assicurati di usare un vero provider di Embedding e reinizia l'ingestion pulendo prima il database.
\end{itemize}

\end{document}
